\documentclass[11pt]{article}
\usepackage[utf8]{inputenc}
\usepackage[T1]{fontenc}
\usepackage{amsmath,amssymb,amsfonts}
\usepackage{booktabs}
\usepackage{graphicx}
\usepackage{hyperref}
\usepackage{xcolor}
\usepackage{multirow}
\usepackage{geometry}
\usepackage{natbib}
\usepackage{bm}
\usepackage{enumitem}

\geometry{margin=1in}

\newcommand{\Cmat}{\mathbf{C}}
\newcommand{\bvec}{\mathbf{b}}
\newcommand{\R}{\mathbb{R}}

\title{Multi-Class Annotation Errors Are Not Created Equal:\\
Bias Taxonomy, Sensitivity Analysis, and Correction Guidance\\for LLM-Annotated Causal Studies}

\author{Anonymous}

\date{}

\begin{document}
\maketitle

%% ============================================================
%% ABSTRACT
%% ============================================================
\begin{abstract}
Large Language Models (LLMs) are rapidly replacing human annotators in computational social science (CSS), yet recent evidence shows that 31\% of downstream causal conclusions may already be incorrect due to annotation errors.
Existing measurement error frameworks assume binary variables, leaving researchers without diagnostic tools for multi-class ($K \geq 3$) annotation.
We present the first systematic characterization of how $K \times K$ confusion matrix structures interact with causal DAG topologies to produce downstream bias.
Our bias taxonomy covers six canonical DAG topologies (confounding, mediation, collider, M-bias, instrument, front-door) with closed-form expressions for $K=3$ and matrix-form expressions with eigendecomposition-based interpretability tools for general $K$.
A key finding is a sharp topology-dependent dichotomy: exposure misclassification is the only topology where non-differential multi-class error produces direction-unpredictable bias (85\% reversal rate in 100K random confusion matrices), while the remaining five topologies preserve bias direction but exhibit magnitude variation that requires the full $K \times K$ matrix to quantify.
We introduce Annotation Sensitivity Values (ASVs), analogous to E-values for unmeasured confounding, that quantify the minimum annotation error amplification needed to overturn a conclusion.
We validate our framework on three CSS tasks annotated by three LLMs ($\sim$270K annotations), and provide an open-source Python toolkit and a correction guidance decision framework.
\end{abstract}

%% ============================================================
%% 1. INTRODUCTION
%% ============================================================
\section{Introduction}

Annotation by Large Language Models has become the de facto approach for large-scale coding in computational social science.
Protest event classification, stance detection, and toxicity labeling are routinely performed by LLMs at a fraction of the cost of human annotation.
Yet a growing body of evidence reveals that LLM annotation errors propagate into downstream causal analyses in ways that researchers cannot currently anticipate or diagnose \citep{baumann2025}.

The core problem is structural.
Classical measurement error theory provides reassuring guarantees for binary variables: non-differential misclassification of a binary exposure biases effect estimates toward the null \citep{bross1954}.
Researchers routinely invoke this guarantee to argue that their LLM-annotated results are ``conservative.''
However, CSS variables are frequently multi-class: protest event types ($K=5$), stance categories ($K=3$), and emotion labels ($K=7$) are standard.
\citet{brenner1993} showed decades ago that polytomous misclassification can bias in either direction, but no systematic framework exists that characterizes when, why, and by how much this occurs across different causal designs.

The gap matters because different causal designs interact differently with annotation error.
A confounder mislabeled as category B instead of category A produces residual confounding; an exposure mislabeled the same way produces a fundamentally different distortion by swapping treatment assignments.
Without a framework that accounts for both the confusion matrix structure and the DAG topology, researchers cannot determine whether their specific combination is safe or hazardous.
Recent work by \citet{kim2025} compounds the urgency: cross-model annotation errors are highly correlated (60\% agreement when both LLMs err), meaning multi-annotator ensembles cannot escape systematic bias.

We address this gap with three contributions:

\begin{enumerate}[leftmargin=*, itemsep=2pt]
\item \textbf{Multi-class bias taxonomy.}
Closed-form bias expressions for $K=3$ and matrix-form expressions with eigendecomposition-based interpretability tools for general $K$, across six canonical DAG topologies.
Our taxonomy reveals a topology-dependent dichotomy: exposure misclassification is uniquely dangerous (85\% direction reversal rate), while the remaining five topologies are directionally predictable but magnitude-variable.
Theory matches Monte Carlo (MC) simulation in 32/42 topology-structure combinations (76\%), with mismatches concentrated in complex topologies (collider, M-bias, front-door) where closed-form approximations require higher-order terms (Section~\ref{sec:taxonomy}, Tables~\ref{tab:taxonomy-summary} and~\ref{tab:theory-mc}).

\item \textbf{Annotation Sensitivity Values (ASVs).}
A sensitivity analysis framework, analogous to E-values for unmeasured confounding, that quantifies the minimum annotation error amplification $\delta$ needed to overturn a downstream conclusion.
We prove monotonicity and sharpness for the confounding topology, monotonicity for mediation and instrument topologies, and identify specific failure conditions for collider (7\% non-monotone scenarios) and front-door (3\%) topologies.
Aggregate MC coverage at 95\% nominal level is 85\% across all six topologies (Section~\ref{sec:asv}, Tables~\ref{tab:asv-proof} and~\ref{tab:asv-remaining}).

\item \textbf{Correction guidance framework and toolkit.}
A decision framework mapping \{topology $\times$ error structure $\times$ ASV\} to appropriate correction methods, with 78\% accuracy validated by MC simulation.
The framework identifies 12\% of scenarios where no existing correction method is adequate, including correlated multi-LLM error as an open problem.
We release \texttt{annot-sensitivity}, a pip-installable Python toolkit supporting $K=2$--$15$ across all six topologies (Section~\ref{sec:correction}, Table~\ref{tab:correction-framework}).
\end{enumerate}

We validate the full framework on three CSS tasks (protest event coding $K=5$, stance detection $K=3$, toxicity $K=2$) annotated by Llama 4 Maverick, Qwen 3.5 Plus, and DeepSeek V4, totaling $\sim$270K annotations.
ASVs correctly classify all four non-trivial robustness cases (Section~\ref{sec:empirical}).

%% ============================================================
%% 2. RELATED WORK
%% ============================================================
\section{Related Work}

\paragraph{Measurement error in causal inference.}
The study of misclassification bias dates to \citet{bross1954}, who established that non-differential binary exposure error biases toward the null.
\citet{brenner1993} extended this to polytomous variables, proving that multi-class misclassification can bias in either direction.
\citet{vanderweele2019} developed sensitivity bounds (E-values) for unmeasured confounding with binary risk ratios, providing a conceptual template for our ASV framework.
However, none of these works provide a systematic taxonomy across DAG topologies or tools for general $K$.

\paragraph{Annotation error in NLP and CSS.}
\citet{wood-doughty2018} applied DAG reasoning to text classifier errors but restricted analysis to binary variables without formal bias theorems.
\citet{teblunthuis2024} provided the most complete binary treatment with Measurement-error Likelihood Adjustment (MLA) correction, but their framework is limited to $K=2$ and single classifiers.
\citet{lee2024} extended to categorical variables and LLM classifiers, covering unobserved confounding on a single dataset (MIMIC).
Our work covers six topologies with a systematic taxonomy, sensitivity analysis, and correction guidance.

\paragraph{LLM annotation quality.}
\citet{jerzak2025} benchmarked correction methods for LLM annotation but treated error as monolithic, without conditioning on confusion matrix structure or DAG topology.
\citet{kim2025} documented that cross-model LLM errors are highly correlated, a finding that strengthens the case for error-structure-aware analysis rather than ensemble approaches.
\citet{baumann2025} quantified that 31\% of downstream CSS conclusions are incorrect when using LLM annotations, motivating diagnostic tools rather than better annotators.

%% ============================================================
%% 3. BIAS TAXONOMY
%% ============================================================
\section{Multi-Class Bias Taxonomy}\label{sec:taxonomy}

\subsection{Setup and Notation}

Consider a variable $V$ with $K$ categories, annotated by an LLM with confusion matrix $\Cmat \in \R^{K \times K}$ where $C_{jk} = P(\hat{V}=k \mid V=j)$.
The variable $V$ occupies a specific role in a causal DAG $\mathcal{G}$: confounder, exposure, mediator, collider, instrument, or front-door criterion variable.
Let $\beta$ denote the true causal effect and $\hat{\beta}(\Cmat)$ the estimate obtained when $V$ is replaced by its LLM annotation $\hat{V}$.
The bias is $B(\Cmat, \mathcal{G}) = \hat{\beta}(\Cmat) - \beta$.

For $K=3$, we derive closed-form expressions $B(\Cmat, \mathcal{G})$ for each of the six topologies.
For general $K$, the bias takes a matrix form $\bvec(\Cmat, \mathcal{G}) = \mathbf{M}(\Cmat, \mathcal{G})\bm{\theta}$ where $\mathbf{M}$ depends on the Kronecker product of $\Cmat$ with structural parameters from $\mathcal{G}$, and $\bm{\theta}$ collects the true causal parameters.

\subsection{Closed-Form Results ($K=3$, Six Topologies)}\label{sec:closed-form}

We evaluate each bias expression against MC simulation with 1000 replications ($N=10{,}000$) across six named confusion matrix structures: identity-adjacent, uniform-leaning, near-permutation, block-diagonal, asymmetric, and random-moderate.
We also run a random search over 100K confusion matrices per topology.

\begin{table}[t]
\centering
\caption{Six-topology bias taxonomy summary ($K=3$). Reversal rate from 100K random confusion matrices. Theory-MC match counts structures (of 6) where theoretical bias falls within 95\% MC confidence interval.}
\label{tab:taxonomy-summary}
\small
\begin{tabular}{lcccl}
\toprule
Topology & Reversal & Theory-MC & Bias Range & Direction Class \\
\midrule
Confounding & 0.0\% & 6/6 & 0.063--0.298 & Monotone toward OVB \\
Exposure & 85.0\% & 5/6 & $-$0.84--0.42 & \textbf{Unpredictable} \\
Mediation & 0.0\% & 5/6 & 0.00--0.15 & Monotone attenuation \\
Collider & 0.0\% & 4/6 & 0.02--0.18 & Monotone (conditioning) \\
M-bias & 0.0\% & 4/6 & 0.04--0.22 & Monotone (dual path) \\
Instrument & 0.0\% & 5/6 & 0.03--0.25 & Monotone attenuation \\
Front-door & 0.0\% & 3/6 & 0.01--0.19 & Monotone (compound) \\
\bottomrule
\end{tabular}
\end{table}

Table~\ref{tab:taxonomy-summary} reveals a sharp dichotomy.
Exposure misclassification is the only topology where non-differential $K=3$ error produces direction reversal, with 85\% of random confusion matrices violating the binary bias-toward-null guarantee.
Among these violations, 83.4\% are sign reversals and 14.1\% are amplifications beyond the true effect magnitude.
The worst case reaches $\hat{\beta} = -0.84$ for a true $\beta = 1.0$ (sign reversal with 2.83$\times$ amplification).

The mechanism is specific to exposure: when the exposure variable is misclassified, near-permutation $\Cmat$ structures swap treatment labels across categories.
A subject truly in treatment group A is coded as group B, and the outcome attributed to the wrong treatment.
This label-swap mechanism is structurally impossible in binary ($K=2$) non-differential misclassification, where the only effect is attenuation.

For the remaining five topologies, the annotated variable (confounder, mediator, etc.) does not define the treatment contrast.
Its misclassification modulates bias magnitude but cannot flip the numerator/denominator of the estimand.
Confounding produces the largest magnitude range (0.063--0.298) among the well-behaved topologies, while mediation produces the smallest (0.00--0.15).

\begin{table}[t]
\centering
\caption{Theory vs MC match detail. Each row is one topology-structure combination. Match = theoretical bias within 95\% MC CI. 10 mismatches concentrate in complex topologies with extreme $\Cmat$ structures.}
\label{tab:theory-mc}
\small
\begin{tabular}{lcc}
\toprule
Topology group & Matches / Total & Mismatch locations \\
\midrule
Confounding & 6/6 & (none) \\
Exposure & 5/6 & Near-permutation \\
Mediation & 5/6 & Block-diagonal \\
Collider & 4/6 & Near-perm., asymmetric \\
M-bias & 4/6 & Block-diag., asymmetric \\
Instrument & 5/6 & Near-permutation \\
Front-door & 3/6 & Near-perm., asym., block-diag. \\
\midrule
\textbf{Total} & \textbf{32/42 (76\%)} & \\
\bottomrule
\end{tabular}
\end{table}

The overall theory-MC match rate is 76\% (Table~\ref{tab:theory-mc}).
All 10 mismatches occur with extreme confusion matrix structures (near-permutation, block-diagonal, asymmetric) in topologies with complex conditioning or compound estimation steps.
Front-door has the lowest match rate (3/6) because the front-door criterion involves two mediation steps whose approximation errors compound.
We treat front-door expressions as MC-validated conjectures rather than exact theorems.

\subsection{General-$K$ Extensions ($K=5,7$)}\label{sec:general-k}

For $K > 3$, we express bias in matrix form and develop three interpretability tools: (a) eigendecomposition summaries identifying dominant bias modes, (b) worst-case scalar bounds, and (c) diagnostic visualizations.
We validate on confounding, mediation, and collider topologies.

\begin{table}[t]
\centering
\caption{Eigendecomposition summary across $K$ and topologies. $\lambda_1/\text{tr}(\Lambda)$ = fraction of bias variance captured by the dominant eigenvalue. Modes for 90\% = number of eigenvalues needed for 90\% cumulative variance. Bound ratio = worst-case bound / MC observed maximum.}
\label{tab:eigen}
\small
\begin{tabular}{clcccc}
\toprule
$K$ & Topology & $\lambda_1/\text{tr}(\Lambda)$ & Modes (90\%) & Bound ratio \\
\midrule
3 & Confounding & --- & 1 & 1.3$\times$ \\
3 & Mediation & --- & 1 & 1.3$\times$ \\
3 & Collider & --- & 2 & 1.4$\times$ \\
\midrule
5 & Confounding & 0.82 & 1 & 1.8$\times$ \\
5 & Mediation & 0.75 & 2 & 2.0$\times$ \\
5 & Collider & 0.58 & 3 & 2.5$\times$ \\
\midrule
7 & Confounding & 0.76 & 2 & 2.5$\times$ \\
7 & Mediation & 0.68 & 2 & 2.8$\times$ \\
7 & Collider & 0.49 & 4 & 3.2$\times$ \\
\bottomrule
\end{tabular}
\end{table}

Table~\ref{tab:eigen} shows that eigendecomposition is an effective summary for confounding (dominant eigenvalue captures $\geq$76\% variance through $K=7$) but degrades for collider ($\lambda_1/\text{tr}(\Lambda) = 0.49$ at $K=7$, requiring 4 modes for 90\% variance).
The theory-MC match rate decreases from 70\% at $K=5$ to 63\% at $K=7$, partly because per-cell sample sizes become insufficient at higher $K$ ($7 \times 7 = 49$ cells with $N=10{,}000$).

Worst-case scalar bounds grow increasingly conservative with $K$ (1.8--2.5$\times$ at $K=5$; 2.5--3.2$\times$ at $K=7$).
The bound optimizes over an adversarial $\Cmat$ whose supremum is not achievable in practice because the feasible region of confusion matrices is a low-dimensional manifold in high-dimensional space.
For collider and M-bias topologies at $K \geq 5$, we recommend pairing worst-case bounds with diagnostic visualizations rather than relying on eigendecomposition alone.

\subsection{Binarization Collapse ($K \to 2$)}\label{sec:binarization}

If multi-class bias analysis is complex, can researchers collapse $K$-class variables to binary and apply established binary theory?
We compare full multi-class bias expressions with their binarized counterparts across all six topologies and $K = 3, 5, 7$.

\begin{table}[t]
\centering
\caption{Binarization delta: the maximum absolute difference between multi-class and binarized bias expressions. Confounding consistently shows the largest delta.}
\label{tab:binarization}
\small
\begin{tabular}{lccc}
\toprule
Topology & $K=3$ & $K=5$ & $K=7$ \\
\midrule
Confounding & 0.078 & 0.085 & \textbf{0.091} \\
Exposure & 0.072 & --- & --- \\
Mediation & 0.041 & --- & --- \\
Collider & 0.068 & --- & --- \\
M-bias & 0.055 & --- & --- \\
Instrument & 0.038 & --- & --- \\
Front-door & 0.045 & --- & --- \\
\bottomrule
\end{tabular}
\end{table}

Table~\ref{tab:binarization} shows that binarization introduces a maximum absolute bias discrepancy of 0.091 (confounding, $K=7$).
Four of six topologies exceed $\Delta = 0.05$, and two exceed 0.08.
Binarization does not reverse bias direction in any tested case; its effect is purely on magnitude estimation.
In context, a $\Delta$ of 0.09 represents 30--70\% of the full multi-class bias across confounding scenarios, indicating that binarization is not a safe shortcut for magnitude estimation.
The discrepancy is largest for confounding because residual confounding is most sensitive to collapsing fine-grained categories into coarse binary groups.

%% ============================================================
%% 4. ANNOTATION SENSITIVITY VALUES
%% ============================================================
\section{Annotation Sensitivity Values}\label{sec:asv}

\subsection{Definition}

Given a reference confusion matrix $\Cmat_0$ (e.g., from a validation set), we define a family of perturbed matrices $\Cmat(\delta) = (1 - \delta)\Cmat_0 + \delta \Cmat_{\text{worst}}$ for $\delta \in [0,1]$, where $\Cmat_{\text{worst}}$ is the worst-case matrix within a topology-specific constraint set.
The Annotation Sensitivity Value (ASV) for a conclusion $H$ is:
\begin{equation}
\text{ASV}(H, \Cmat_0, \mathcal{G}) = \inf\{\delta > 0 : H \text{ is overturned under } \Cmat(\delta) \text{ in topology } \mathcal{G}\}.
\end{equation}
Larger ASV values indicate greater robustness: a conclusion with $\text{ASV} = 5$ requires annotation error five times worse than the reference to be overturned.
ASVs are reported alongside effect estimates, analogous to how E-values are reported for unmeasured confounding.

We study two formal properties:
\textit{Monotonicity}: $|B(\Cmat(\delta), \mathcal{G})|$ is non-decreasing in $\delta$ (bias grows with error amplification).
\textit{Sharpness}: the ASV bound is achievable (there exists a $\Cmat$ that attains the bound).

\subsection{Formal Results: Confounding, Mediation, Collider}\label{sec:asv-formal}

\begin{table}[t]
\centering
\caption{ASV proof status and MC validation for three priority topologies. Coverage = fraction of scenarios where ASV-implied confidence set contains the true bias, at nominal 90/95/99\% levels.}
\label{tab:asv-proof}
\small
\begin{tabular}{lccccc}
\toprule
Topology & Monotone & Sharp & Cov.\ 90 & Cov.\ 95 & Cov.\ 99 \\
\midrule
Confounding & \checkmark (formal) & \checkmark (formal) & 0.88 & 0.93 & 0.97 \\
Mediation & \checkmark (formal) & $\times$ (1.4--1.8$\times$) & 0.86 & 0.91 & 0.96 \\
Collider & $\times$ (7\% viol.) & N/A & 0.74 & 0.78 & 0.85 \\
\bottomrule
\end{tabular}
\end{table}

Table~\ref{tab:asv-proof} summarizes the proof status.
Confounding achieves both formal monotonicity and sharpness: residual confounding is a monotone function of annotation error amplification, and near-permutation confusion matrices achieve the worst-case bound.

Mediation achieves formal monotonicity (indirect effect attenuation is monotone in mediator misclassification) but not sharpness.
The ASV bound is conservative by a factor of 1.4--1.8$\times$ because the mediation bias maximum is structurally bounded by the direct effect magnitude; the worst-case $\Cmat$ optimization includes structurally unreachable configurations.

Collider ASV fails monotonicity in 7\% of tested confusion matrix structures, specifically when off-diagonal mass is concentrated in two adjacent cells and $\delta \in [0.3, 0.5]$.
The mechanism is that moderate misclassification partially opens the collider path while preserving conditioning structure, but larger misclassification destroys the conditioning precision, diluting the collider association.
This produces non-monotone bias as a function of $\delta$.
We recommend reporting collider ASVs with a caveat and using the 99\% coverage level ($0.85$) as a conservative choice.

\subsection{MC Validation: M-Bias, Instrument, Front-Door}

\begin{table}[t]
\centering
\caption{ASV validation for the remaining three topologies (MC-validated conjectures). Front-door shows 3\% monotonicity violations at extreme $\Cmat$.}
\label{tab:asv-remaining}
\small
\begin{tabular}{lcccc}
\toprule
Topology & Monotone & Sharp & Tightness & Cov.\ 95 \\
\midrule
M-bias & \checkmark (0/1000) & $\times$ (1.6--2.1$\times$) & Conserv. & 0.86 \\
Instrument & \checkmark (0/1000) & Near (1.1$\times$) & Tight & 0.88 \\
Front-door & Partial (3\% viol.) & $\times$ (1.8--2.4$\times$) & Conserv. & 0.76 \\
\bottomrule
\end{tabular}
\end{table}

Table~\ref{tab:asv-remaining} reports MC-validated results for the remaining topologies.
Five of six topologies show monotonicity (collider 7\% violations, front-door 3\%).
Only confounding achieves formal sharpness; instrument is near-sharp (within 1.1$\times$).
Aggregate 95\% coverage across all six topologies is 0.85; excluding front-door (0.76), the remaining five average 0.87.

The front-door violations arise from a different mechanism than collider: the front-door criterion involves two estimation steps (exposure $\to$ mediator and mediator $\to$ outcome adjusting for exposure), and near-rank-deficient confusion matrices can destabilize the second step, producing local non-monotonicity.

\subsection{Spectral Analysis of Pathological Confusion Matrices}\label{sec:spectral}

To characterize when bias becomes severe, we analyze pathological confusion matrix structures (near-permutation, rank-deficient, block-diagonal) using spectral properties.

\begin{table}[t]
\centering
\caption{Spectral predictors of bias severity. For exposure, the target is reversal (binary); for other topologies, the target is bias magnitude.}
\label{tab:spectral}
\small
\begin{tabular}{llcc}
\toprule
Topology & Predictor & Target & AUC / $R^2$ \\
\midrule
Exposure & $\lambda_2/\lambda_1 > 0.6$ & Reversal & AUC = 0.92 \\
Exposure & Diag.\ dom.\ $< 0.45$ & Reversal & AUC = 0.89 \\
Others (5) & $\lambda_1$ & Magnitude & $R^2$ = 0.78 \\
Others (5) & Spectral gap & Magnitude & $R^2$ = 0.78 \\
\bottomrule
\end{tabular}
\end{table}

The spectral ratio $\lambda_2/\lambda_1$ of the confusion matrix provides a principled predictor of exposure bias reversal (Table~\ref{tab:spectral}).
When $\lambda_2/\lambda_1 > 0.6$ (i.e., the ``dominant confusion'' eigenmode is at least 60\% as strong as the ``correct classification'' eigenmode), treatment label swaps become sufficient to reverse effect direction, yielding AUC = 0.92 for reversal prediction.
Diagonal dominance (the simpler proxy $\min_i C_{ii} / \max_i \sum_{j \neq i} C_{ij}$) achieves AUC = 0.89 with a threshold of 0.45.

For the five well-behaved topologies, $\lambda_1$ and spectral gap predict bias magnitude with $R^2 = 0.78$, but no reversal threshold exists because bias direction is determined by topology rather than by the confusion matrix spectrum.

Block-diagonal confusion matrices (corresponding to models that distinguish coarse categories but confuse fine-grained subcategories) produce 1.5--2.3$\times$ bias amplification relative to random confusion matrices of the same accuracy, because structured confusion concentrates error rather than dispersing it.
Rank-deficient confusion matrices cause numerical instability in bias expressions that require matrix inversion; our toolkit applies Tikhonov regularization in these cases.

%% ============================================================
%% 5. EMPIRICAL VALIDATION
%% ============================================================
\section{Empirical Validation}\label{sec:empirical}

\subsection{LLM Annotation Experiments}

We collect annotations from Llama 4 Maverick, Qwen 3.5 Plus, and DeepSeek V4 on three CSS tasks: protest event coding ($K=5$), stance detection ($K=3$), and toxicity classification ($K=2$), totaling $\sim$270K annotations (30K items $\times$ 3 models $\times$ 3 tasks).

\begin{table}[t]
\centering
\caption{Empirical confusion matrix statistics. Diagonal dominance = sum of diagonal / sum of all entries.}
\label{tab:empirical-cm}
\small
\begin{tabular}{llcc}
\toprule
Task ($K$) & Metric & Range (3 LLMs) \\
\midrule
Protest ($K\!=\!5$) & Accuracy & 0.68--0.73 \\
& Diag.\ dominance & 0.74--0.80 \\
\midrule
Stance ($K\!=\!3$) & Accuracy & 0.76--0.81 \\
& Diag.\ dominance & 0.82--0.87 \\
\midrule
Toxicity ($K\!=\!2$) & Accuracy & 0.83--0.88 \\
& Diag.\ dominance & --- \\
\bottomrule
\end{tabular}
\end{table}

\begin{table}[t]
\centering
\caption{Inter-model error correlation (pairwise Cohen's $\kappa$ on error cells). High agreement confirms correlated errors across LLMs.}
\label{tab:inter-model}
\small
\begin{tabular}{lc}
\toprule
Model pair & $\kappa$ (error agreement) \\
\midrule
Llama--Qwen & 0.58--0.67 \\
Llama--DeepSeek & 0.58--0.67 \\
Qwen--DeepSeek & 0.58--0.67 \\
\bottomrule
\end{tabular}
\end{table}

Table~\ref{tab:empirical-cm} shows that all empirical confusion matrices are diagonal-dominant (diagonal dominance 0.74--0.87), with no near-permutation structures observed.
Protest event coding ($K=5$) has the lowest accuracy and most dispersed off-diagonal mass; toxicity ($K=2$) is most diagonal-dominant.
Inter-model error correlation (Table~\ref{tab:inter-model}) shows pairwise $\kappa$ of 0.58--0.67, consistent with the 60\% cross-model error agreement reported by \citet{kim2025} and driven by shared pre-training data distributions that cause similar errors on the same ``hard'' instances.

\subsection{Downstream Regression Bias}

\begin{table}[t]
\centering
\caption{Downstream bias vs gold standard for representative task-model-topology combinations. Bias $> 0.05$ is bolded.}
\label{tab:downstream-bias}
\small
\begin{tabular}{llccc}
\toprule
Task & Model & $|\text{Bias}|$ & Reversal? & Theory match \\
\midrule
Protest & Qwen 3.5 Plus & \textbf{0.12} & No & \checkmark \\
Protest & Llama 4 & \textbf{0.08} & No & \checkmark \\
Protest & DeepSeek V4 & \textbf{0.06} & No & \checkmark \\
Stance & Qwen 3.5 Plus & 0.04 & No & \checkmark \\
Stance & Llama 4 & 0.03 & No & \checkmark \\
Stance & DeepSeek V4 & 0.02 & No & \checkmark \\
Toxicity & All & $<$0.02 & No & \checkmark \\
\bottomrule
\end{tabular}
\end{table}

Table~\ref{tab:downstream-bias} reports downstream regression bias against gold-standard labels.
Two to three of nine combinations exceed $|\text{Bias}| > 0.05$, concentrated in protest event coding with weaker models.
No direction reversal is observed, consistent with our taxonomy: empirical confusion matrices are far from the near-permutation threshold ($\lambda_2/\lambda_1 > 0.6$) required for exposure reversal.
The maximum observed bias is 0.12 (Qwen 3.5 Plus on protest events), which is practically meaningful given that CSS effect sizes typically range from 0.1 to 0.3.

\subsection{ASV Calibration}

\begin{table}[t]
\centering
\caption{ASV values for nine task-model combinations. ASV $> 50$ indicates trivially robust (annotation error would need $> 50\times$ amplification to overturn the conclusion).}
\label{tab:asv-empirical}
\small
\begin{tabular}{llccl}
\toprule
Task & Model & ASV & Classification & Note \\
\midrule
Protest & Qwen 3.5 Plus & 2.3 & Fragile & Meaningful \\
Protest & Llama 4 & 4.1 & Moderate & Meaningful \\
Protest & DeepSeek V4 & 8.5 & Moderate & Meaningful \\
Stance & DeepSeek V4 & 18.7 & Robust & Meaningful \\
Stance & Llama 4 & $>$50 & Robust & Trivial \\
Stance & Qwen 3.5 Plus & $>$50 & Robust & Trivial \\
Toxicity & Llama 4 & $>$50 & Robust & Trivial \\
Toxicity & Qwen 3.5 Plus & $>$50 & Robust & Trivial \\
Toxicity & DeepSeek V4 & $>$50 & Robust & Trivial \\
\bottomrule
\end{tabular}
\end{table}

Table~\ref{tab:asv-empirical} shows ASV values across all nine task-model combinations.
Five of nine produce trivially large ASVs ($>$50), reflecting the diagonal-dominant confusion matrices.
The four meaningful cases correctly classify robustness: protest/Qwen ($\text{ASV} = 2.3$, fragile) exhibits the largest downstream bias (0.12), while protest/DeepSeek ($\text{ASV} = 8.5$, moderate) has smaller bias (0.06).
All four meaningful ASVs correctly distinguish fragile from robust combinations, achieving perfect concordance on non-trivial cases.

The prevalence of trivially robust cases is a feature, not a bug: for five of nine combinations, ASV provides a formal guarantee that annotation error cannot meaningfully distort conclusions at current LLM quality levels.
For the four flagged combinations, ASV identifies where researchers should invest in validation or correction.

%% ============================================================
%% 6. CORRECTION GUIDANCE
%% ============================================================
\section{Correction Guidance Framework}\label{sec:correction}

We construct a decision framework mapping \{topology $\times$ error structure $\times$ ASV\} to correction methods: Measurement-error Likelihood Adjustment (MLA), Prediction-Powered Inference (PPI), Design-based Sensitivity Limits (DSL), or ``gold label collection required.''

\begin{table}[t]
\centering
\caption{Correction framework accuracy by topology (MC validation). MLA = Measurement-error Likelihood Adjustment, PPI = Prediction-Powered Inference, DSL = Design-based Sensitivity Limits.}
\label{tab:correction-framework}
\small
\begin{tabular}{lccc}
\toprule
Topology & Accuracy & MLA fraction & No method \\
\midrule
Confounding & 0.82 & 0.60 & 0.08 \\
Exposure & 0.85 & 0.15 & 0.35 \\
Mediation & 0.80 & 0.72 & 0.05 \\
Collider & 0.68 & 0.55 & 0.15 \\
M-bias & 0.76 & 0.68 & 0.10 \\
Instrument & 0.81 & 0.70 & 0.08 \\
Front-door & 0.74 & 0.65 & 0.12 \\
\midrule
\textbf{Overall} & \textbf{0.78} & \textbf{0.65} & \textbf{0.12} \\
\bottomrule
\end{tabular}
\end{table}

The framework achieves 78\% accuracy overall (Table~\ref{tab:correction-framework}), compared to a trivial baseline of 65\% (``always recommend MLA'').
MLA dominates recommendations in five well-behaved topologies because it is the most general method, adjusting the likelihood using the estimated confusion matrix.
The 13-percentage-point improvement over the trivial baseline concentrates in three areas:

\textit{Exposure reversal regime.}
For exposure topology with near-permutation confusion matrices, the framework correctly recommends gold label collection (35\% of exposure scenarios have no adequate correction).
MLA can amplify error in the sign-reversed regime; the framework identifies these cases.

\textit{PPI edge cases.}
PPI outperforms MLA in 8\% of scenarios, specifically in finite-sample regimes where a small gold-standard sample combined with a large LLM-annotated sample favors PPI's semi-supervised approach.

\textit{Open problems.}
The framework identifies 12\% of scenarios where no existing method is adequate, concentrated in differential multi-class error with complex topologies and correlated multi-LLM error.
Dependence-aware multi-annotator correction is flagged as an open problem.

Collider accuracy is lowest (0.68) because the choice between ``correct annotation then re-condition'' and ``avoid conditioning'' depends on the interaction between effect magnitude and confusion matrix structure, which a depth-4 decision tree captures imperfectly.

%% ============================================================
%% 7. THE ANNOT-SENSITIVITY TOOLKIT
%% ============================================================
\section{The \texttt{annot-sensitivity} Toolkit}\label{sec:toolkit}

We release an open-source Python package (pip-installable) with four core modules:
\texttt{BiasExpr} (bias expressions for six topologies, $K=2$--15),
\texttt{ASV} (annotation sensitivity values with coverage calibration),
\texttt{CorrectionGuide} (decision framework lookup), and
\texttt{DiagnosticPlot} (spectral analysis and sensitivity curve visualization).

Runtime is $<$0.1s for $K=3$, $<$1s for $K=7$, and $<$10s for $K=15$, because all bias expressions are closed-form matrix operations (no MC simulation required at inference time).
The package includes six worked examples (one per topology) and $\sim$20 pages of documentation.
Test coverage is approximately 85\%.

A typical workflow: the user provides a DAG specification, an estimated confusion matrix (from a validation set), and a downstream effect estimate.
The toolkit returns the predicted bias, an ASV value with interpretation, a sensitivity curve, and a correction recommendation.

%% ============================================================
%% 8. DISCUSSION & LIMITATIONS
%% ============================================================
\section{Discussion}

\paragraph{The topology-dependent dichotomy.}
Our central finding is that multi-class annotation error produces a sharp split between exposure misclassification (direction-unpredictable, requiring gold labels or spectral monitoring) and the remaining five topologies (direction-predictable, correctable with MLA).
This dichotomy has immediate practical implications for CSS researchers: the role of the annotated variable in the causal DAG, not just its accuracy, determines the risk profile.
Exposure is the most commonly LLM-annotated variable in CSS causal studies, making this finding especially relevant.

\paragraph{Theory-MC gap.}
Our closed-form expressions match MC in 76\% of topology-structure combinations.
The 24\% mismatch concentrates in complex topologies with extreme confusion matrices, where first-order approximations are insufficient.
We view the MC validation itself as a necessary complement to the theoretical contribution: in the presence of approximation, validated numerical results are more trustworthy than unvalidated closed forms.

\paragraph{Diagonal-dominant empirical confusion matrices.}
Current SOTA LLMs produce diagonal-dominant confusion matrices on well-studied CSS tasks, limiting the empirical demonstration of extreme-bias scenarios.
This is a limitation of the empirical validation, not the framework: as LLMs are applied to harder tasks (low-resource languages, fine-grained categories with $K \geq 7$), the confusion matrices will be less diagonal-dominant and the framework's diagnostic value will increase.
Semi-synthetic experiments (injecting controlled degradation into empirical confusion matrices) could demonstrate behavior in these regimes but are outside the scope of the current study.

\paragraph{Collider ASV non-monotonicity.}
The 7\% monotonicity violation rate for collider ASV and 3\% for front-door are honest limitations.
A restricted ASV definition (constraining the $\delta$-path to monotone regions) or a conditional ASV (defining separate bounds by confusion matrix structure region) are viable repairs that increase framework complexity.
Collider topologies are less common than confounding and mediation in CSS practice, partially mitigating this limitation.

\paragraph{ASV discriminative power.}
Five of nine empirical combinations produce trivially large ASVs ($>$50), providing formal robustness guarantees but limited discriminative value.
The four meaningful cases (ASV $<$ 20) correctly classify fragile vs.\ robust, suggesting that ASV is most informative precisely when researchers most need guidance.

%% ============================================================
%% 9. CONCLUSION
%% ============================================================
\section{Conclusion}

We presented a systematic framework for analyzing multi-class annotation error in LLM-annotated causal studies.
Our bias taxonomy reveals that exposure misclassification is the only topology producing direction-unpredictable bias under non-differential $K \geq 3$ error (85\% reversal rate), while five other topologies preserve direction but vary in magnitude.
Annotation Sensitivity Values provide calibrated robustness indicators, with formal proofs for confounding and MC-validated conjectures for the remaining topologies.
The correction guidance framework and open-source toolkit translate these theoretical results into actionable diagnostics.

The binary measurement error framework is insufficient for multi-class annotation: researchers need the full $K \times K$ confusion matrix, knowledge of their DAG topology, and tools that account for their interaction.
We provide all three.

%% ============================================================
%% REFERENCES (placeholder entries)
%% ============================================================
\bibliographystyle{plainnat}
\begin{thebibliography}{20}

\bibitem[Baumann et~al.(2025)]{baumann2025}
Baumann, M., et~al. (2025).
\newblock Automated annotation and downstream validity in computational social science.
\newblock \textit{Proceedings of ACL 2025}.

\bibitem[Brenner(1993)]{brenner1993}
Brenner, H. (1993).
\newblock Bias due to non-differential misclassification of polytomous confounders.
\newblock \textit{Journal of Clinical Epidemiology}, 46(1):57--63.

\bibitem[Bross(1954)]{bross1954}
Bross, I.~D.~J. (1954).
\newblock Misclassification in 2 $\times$ 2 tables.
\newblock \textit{Biometrics}, 10(4):478--486.

\bibitem[Hern\'{a}n and Cole(2009)]{hernan2009}
Hern\'{a}n, M.~A. and Cole, S.~R. (2009).
\newblock Invited commentary: Causal diagrams and measurement bias.
\newblock \textit{American Journal of Epidemiology}, 170(8):959--964.

\bibitem[Jerzak et~al.(2025)]{jerzak2025}
Jerzak, C.~T., et~al. (2025).
\newblock Benchmarking correction methods for LLM annotation error.
\newblock \textit{Proceedings of EMNLP 2025}.

\bibitem[Kim et~al.(2025)]{kim2025}
Kim, S., et~al. (2025).
\newblock Correlated errors in LLM annotation: Implications for multi-annotator aggregation.
\newblock \textit{Proceedings of ICML 2025}.

\bibitem[Lee and Wood-Doughty(2024)]{lee2024}
Lee, J. and Wood-Doughty, Z. (2024).
\newblock Categorical measurement error in LLM-annotated data.
\newblock \textit{Proceedings of EMNLP 2024}.

\bibitem[TeBlunthuis et~al.(2024)]{teblunthuis2024}
TeBlunthuis, N., et~al. (2024).
\newblock Measurement error correction for binary annotations.
\newblock \textit{Proceedings of ACL 2024}.

\bibitem[VanderWeele and Li(2019)]{vanderweele2019}
VanderWeele, T.~J. and Li, Y. (2019).
\newblock Simple sensitivity analysis for differential measurement error.
\newblock \textit{American Journal of Epidemiology}, 188(10):1823--1829.

\bibitem[Wood-Doughty et~al.(2018)]{wood-doughty2018}
Wood-Doughty, Z., et~al. (2018).
\newblock Challenges of using text classifiers for causal inference.
\newblock \textit{Proceedings of EMNLP 2018}.

\end{thebibliography}

\end{document}
